% Generated from project outputs. Do not edit by hand.
\begin{table}[!htbp]
\centering
\begin{threeparttable}
\caption{Labor income by subsample: 2021-2024}
\label{tab:descriptive-2021-2024-subsamples-y1}
\small
\setlength{\tabcolsep}{4pt}
\renewcommand{\arraystretch}{1.15}
\begin{tabularx}{\linewidth}{@{}Xrrrrr@{}}
\toprule
Sample & N & Mean t & SD t & Mean t+1 & SD t+1 \\
\midrule
General & 276 & 2,817.9 & 1,499.6 & 3,130.2 & 1,556.9 \\
Men & 276 & 3,185.9 & 1,624.5 & 3,564.0 & 1,660.9 \\
Women & 276 & 2,231.9 & 1,540.6 & 2,493.5 & 1,614.8 \\
Urban & 276 & 3,041.8 & 1,557.9 & 3,332.9 & 1,596.8 \\
Rural & 276 & 2,488.8 & 1,633.1 & 2,781.8 & 1,745.8 \\
Indigenous & 276 & 2,477.3 & 1,771.6 & 2,924.3 & 1,814.9 \\
Non-Indigenous & 184 & 2,973.8 & 1,566.2 & 3,184.4 & 1,506.8 \\
Bottom 99 percent & 276 & 2,641.3 & 1,285.2 & 2,934.7 & 1,305.7 \\
Bottom 90 percent & 276 & 2,179.7 & 884.4 & 2,410.6 & 866.6 \\
Bottom 50 percent & 276 & 1,212.0 & 479.7 & 1,326.6 & 450.4 \\
\bottomrule
\end{tabularx}
\begin{tablenotes}[flushleft]
\footnotesize
\item Monthly income in quetzales. Unweighted descriptive statistics of survey-weighted cohort means. N counts cohort-transition rows, not individuals or independent cohorts. SD is dispersion across cohort means, not the standard error of a mean. Both incomes must be observed. Subsamples overlap; bottom-income groups are cumulative.
\end{tablenotes}
\end{threeparttable}
\end{table}
